\section{F. Statement and literature registers}
The accompanying \texttt{RESEARCH\_SA06\_results.csv} is the statement
register. The report begins with one row for every task and every Lead,
including the unresolved rows. Each verification event is recorded in
\texttt{CHECK\_EVENTS.tsv}; the dated \texttt{PROGRESS\_LEDGER.md} records
actual changes of status or reductions of a gate. We do not count
typesetting or repackaging as mathematical movement.

The bibliography marks sources as opened or not obtained. The report's
literature register names the theorem used, the nearest theorem where the
precise requested claim was not obtained, and unsuccessful retrievals.
All cited unpublished manuscripts are identified as such. Numerical
certificates are finite checks of the examples stated, not experimental
proofs of unrestricted assertions.

\paragraph{Final stopping decision.}
We stopped mathematical research at the scheduled 20:50 UTC review on
19 November 2026, after two consecutive reviews without a change in an
open main-row status or a proved reduction to a strictly smaller main gate.
The mathematical programme remains \textbf{INCOMPLETE}.
The final report and dated ledger retain all 41 Part and Lead rows,
including every unresolved gate. Subsequent compilation, visual inspection
and archiving implement delivery and do not constitute mathematical movement.
Dated intermediate notes are retained with their original scope and are
superseded only where a later certified statement explicitly says so.

\paragraph{F.2a (Proved correction): two normalizations in the supplied Tate survey.}
In the notation of \citet[Theorem 1.6 and its proof, p.~3]{rootMilneTate2007},
the sufficient hypothesis is $S^d(X\times X,\ell)$ for $\dim X=d$.
The printed $S^{2d}$ does not match the cohomology group in that proof:
$S^r$ is defined using $H^{2r}(-,\mathbb Q_\ell(r))$.
To check the corrected implication, extend coefficients to split
Frobenius. Poincar\'e duality and K\"unneth identify the tensor of
a generalized eigenvalue-$a$ block $V_a$ and its reciprocal block
with $\operatorname{End}(V_a)$ in
$H^{2d}(X\times X,\mathbb Q_\ell(d))$.
Frobenius acts on this endomorphism space by conjugation. If its
unipotent Jordan factor $U$ were nontrivial, conjugation by $U$
would be a nontrivial unipotent operator: an endomorphism commuting
with every endomorphism is scalar, and a scalar unipotent operator
is the identity. The corrected $S^d$ hypothesis rules this out,
and all Frobenius blocks are semisimple. We do not extract an
unconditional semisimplicity theorem from the printed top-degree
index.

For a $p^{n_1}$-Frobenius $\pi_1$ and a $p^{n_2}$-Frobenius
$\pi_2$, common finite-field extension gives the crossed relation
\[
                    \pi_1^{n_2N}=\pi_2^{n_1N}.
\]
The survey's preceding discussion on p.~2 has these powers; its
definition on p.~4 interchanges them incorrectly. For example,
on $H^2(\mathbb P^1)$ the Frobenius scalars over $\mathbb F_p$
and $\mathbb F_{p^2}$ are $p$ and $p^2$. Crossed powers agree,
whereas the uncrossed equality would say $p^N=p^{4N}$ for
$N>0$, which is impossible. These are normalization checks,
not assertions that the survey's correctly normalized inputs fail.

\paragraph{F.2b (Known, with the supplied version corrected): numerical positivity cones.}
The published \citet[Theorem 2.1 and Corollary 2.2,
pp.~1799--1801]{rootDELV2011} give, on an $n$-fold power of a
CM elliptic curve over $\mathbb C$, pseudoeffective classes equal to strongly
positive forms and nef classes equal to weakly positive forms.
The resulting strict difference has range $2\le k\le n-2$;
the supplied early TeX's Corollary 2.2 instead prints the
incorrect upper endpoint $n$. The published proof also correctly
uses the \emph{closure} of the cone of rank-one forms from
CM-lattice directions. Indeed the CM fraction field is dense
in $\mathbb C$; approximate the coefficients of a one-form,
clear denominators, and rescale its rank-one positive form.
Each approximating form is a positive multiple of the pullback
of the origin under a homomorphism $E^n\to E$.
Closedness gives the asserted pseudoeffective limit.
This verifies the step without the early version's stronger
finite-combination wording. These are numerical cycle cones,
distinct from the character cones in E.1.
